\section{技术基础与统一分析框架}
\label{sec:foundations}

\subsection{从参数生成到检索条件生成}

设用户输入为 $x$，外部知识库为 $\mathcal D$。检索器根据查询表示从 $\mathcal D$ 中返回得分最高的 $k$ 个文档，可写为
\begin{equation}
 E=\operatorname{TopK}_{d\in\mathcal D}\,s_{\phi}(x,d),
 \label{eq:retrieval}
\end{equation}
其中 $s_{\phi}$ 是查询与文档的相关性函数，$E$ 是提供给生成器的证据集合。最基本的生成过程为 $y\sim p_{\theta}(y\mid x,E)$。在原始 RAG 的序列公式中，模型还可以对多个检索文档的预测进行边缘化，而不是简单地把所有文本拼接为唯一上下文~\citep{lewis2020rag}。无论具体实现如何，这一抽象都揭示了两个知识来源：参数 $\theta$ 中隐含的模型知识，以及推理时显式提供的 $E$。

式~\eqref{eq:retrieval} 看似简单，却包含四个容易失效的环节。首先，$x$ 未必适合作为查询，口语化、指代或多跳问题常需改写；其次，相关性分数高不代表文档包含回答问题所需的全部事实；再次，把证据放入上下文不代表模型会正确使用；最后，最终答案的流畅性也不能证明其受到 $E$ 支持。因此，一个完整的故障链可以写为“需求误判 $\rightarrow$ 查询偏差 $\rightarrow$ 检索遗漏或噪声 $\rightarrow$ 上下文利用失败 $\rightarrow$ 无依据生成”。后续方法虽名称各异，实质上都在这一链条的某个位置增加反馈。

\subsection{自适应闭环：状态、行动与反馈}

为统一比较不同方法，本文把第 $t$ 轮的系统状态概念性表示为
\begin{equation}
 s_t=(x,G_t,E_t,b_t),
 \label{eq:state}
\end{equation}
其中 $G_t$ 是已经生成的内容或中间推理，$E_t$ 是当前证据，$b_t$ 是剩余的检索、Token 或时间预算。系统据此选择行动 $a_t$：直接回答、改写查询、检索新证据、过滤或补充现有证据，以及停止并输出答案。行动后状态和证据更新，再进入下一轮判断。该表示是本文为比较方法所作的概念抽象，并非所有论文都显式实现了这一状态结构。

这一视角将常见方法分成四类。\textbf{触发决策}回答“现在是否需要外部知识”，典型信号包括模型置信度、问题复杂度和学习到的检索标记；\textbf{查询决策}回答“下一查应该寻找什么”，包括改写、分解和利用中间推理生成查询；\textbf{证据决策}回答“返回结果是否可用、如何修复”，涉及相关性评价、去噪和替代知识源；\textbf{停止决策}回答“新增检索是否仍值得”，可由固定轮数、答案形成或证据充分性控制。前三类研究较丰富，第四类却常被简化为预设迭代次数。

自适应流程的收益必须与开销共同考虑。本文用下式表达概念性目标：
\begin{equation}
 \max_{\pi}\; \mathbb E_{\pi}\left[Q(y,E)-\lambda_r C_r-\lambda_l C_l-\lambda_u C_u\right],
 \label{eq:objective}
\end{equation}
其中 $\pi$ 是选择行动的策略，$Q$ 表示答案质量及其证据支持程度，$C_r$、$C_l$ 分别表示检索和生成/延迟成本，$C_u$ 表示无依据或错误回答的风险，$\lambda$ 系数体现应用偏好。该式不用于复现某篇论文，也不意味着不同成本总能精确换成同一数值；它的作用是说明，“始终检索”和“始终迭代”一般不是天然最优策略。对医疗、法律等高风险应用，应提高错误风险权重；对交互式闲聊，延迟与额外调用可能更受关注。

由此，自适应 RAG 的核心并不是把系统画得越来越复杂，而是用尽可能少且可靠的行动弥补信息缺口。一个简单分类器如果在明确场景下能以低成本作出正确路由，就可能比多层自我反思更有实用价值；反之，在知识缺口随推理逐步暴露的任务中，单次路由再准确也无法替代后续反馈。
